% Section 3, printed pp. 22--26.
\section{Relative effective divisors}\label{kps-sec-divisors}

\begin{definition}[Effective Cartier divisor]\label{kps-def-effective-divisor}
\dcref{3.1}\source{kleiman-picard:page-0022}
A closed subscheme $D\subset X$ is an effective Cartier divisor when its ideal
$\mathcal I$ is invertible. For any $\mathcal O_X$-module $\mathcal F$ put
$\mathcal F(nD)=\mathcal F\ox\mathcal I^{-n}$. The inclusion
$\mathcal O_X\hookrightarrow\mathcal O_X(D)$ is a regular section. For an
invertible $\mathcal L$, regular sections are those giving injections
$\mathcal L^{-1}\hookrightarrow\mathcal O_X$; $|\mathcal L|$ is the set of
effective divisors with associated sheaf isomorphic to $\mathcal L$.
\end{definition}

\begin{remark}[Complete linear system from regular sections]\label{kps-ex-linear-system}
\source{kleiman-picard:page-0022}
There is a canonical isomorphism
\[H^0(X,\mathcal L)_{\rm reg}/H^0(X,\mathcal O_X^*)\risom|\mathcal L|.\]
\end{remark}

\begin{definition}[Relative effective divisor]\label{kps-def-relative-divisor}
\dcref{3.3}\source{kleiman-picard:page-0022}
A relative effective divisor on $X/S$ is an effective Cartier divisor that is
$S$-flat.
\end{definition}

\begin{lemma}[Local criterion for relative divisors]\label{kps-lem-relative-criterion}
\dcref{3.4}\source{kleiman-picard:page-0022}
For $D\subset X$ and $x\in D$ over $s\in S$, the following are equivalent:
(i) $D$ is a relative effective divisor near $x$; (ii) $X$ and $D$ are
$S$-flat at $x$ and $D_s$ is an effective divisor on $X_s$ at $x$; (iii) $X$
is $S$-flat at $x$ and the ideal of $D$ is generated by an element whose image
is a nonzerodivisor on $\mathcal O_{X_s,x}$.
\end{lemma}
\begin{proof}
Write $A=\mathcal O_{S,s}$, $B=\mathcal O_{X,x}$, $C=\mathcal O_{D,x}$,
and $k=k(s)$. From $0\to B\xrightarrow bB\to C\to0$, flatness of $C$
forces multiplication by $b$ on $B\otimes_Ak$ to be injective; Nakayama and
the local flatness criterion then give (i)$\Rightarrow$(ii). Under (ii), lift a
generator of the fiber ideal and use injectivity of $I\otimes k\to B\otimes k$
and Nakayama to obtain (iii). Under (iii), the Tor sequence shows $C$ flat;
then $I$ is flat, and the exact sequence $0\to K\to B\to I\to0$ after tensoring
with $k$ gives $K\otimes k=0$, hence $K=0$. Thus $b$ is regular and (i) follows.
\end{proof}

\begin{remark}[Sums]\label{kps-ex-sums}
\source{kleiman-picard:page-0023}
The sum of two relative effective divisors is relative effective.
\end{remark}

\begin{definition}[Divisor functor]\label{kps-def-divisor-functor}
\dcref{3.6}\source{kleiman-picard:page-0023}
\[\Div_{X/S}(T)=\{\text{relative effective divisors on }X_T/T\}.\]
Pullback preserves this property because the ideal remains invertible and
$T$-flatness is preserved by base change.
\end{definition}

\begin{theorem}[Representability of relative divisors]\label{kps-thm-divisor-representability}
\dcref{3.7}\source{kleiman-picard:page-0023}
If $X/S$ is projective and flat, $\Div_{X/S}$ is represented by an open
subscheme of $\Hilb_{X/S}$.
\end{theorem}
\begin{proof}
Let $W\subset X\times_S\Hilb_{X/S}$ be the universal subscheme. The locus
$V\subset W$ where $W$ is an effective divisor is open. Its complement has
closed image under the proper projection to the Hilbert scheme; the complement
$U$ of that image is therefore open, and $W_U$ is relative effective. A
relative divisor over $T$ gives a unique map $T\to\Hilb_{X/S}$; each fiber is
a divisor, so the local criterion shows its image lies in $U$. Hence $U$ has
the required universal property.
\end{proof}

\begin{remark}[Curves and symmetric products]\label{kps-ex-curve-divisors}
\source{kleiman-picard:page-0024}
For a flat, locally projective family of pure curves, degree-$m$ divisors are
open-and-closed finite-type pieces $\Div^m_{X/S}$, $\Div^1_{X/S}$ is the
smooth locus, and the sum map from the $m$-fold fiber product of the smooth
locus is natural.\end{remark}
\begin{remark}[Symmetric product]\label{kps-rmk-symmetric-product}
\dcref{3.9}\source{kleiman-picard:page-0024}
The sum map factors through the symmetric product and induces
$X_0^{(m)}\risom\Div^m_{X_0/S}$.
\end{remark}

\begin{definition}[The cohomology module $Q$]\label{kps-def-Q}
\source{kleiman-picard:page-0024}
For proper $f$ and coherent $S$-flat $\mathcal F$, there is a coherent
$\mathcal O_S$-module $Q$ characterized by
\[\Hom(Q,\mathcal N)\risom f_*(\mathcal F\ox f^*\mathcal N)\]
functorially in quasi-coherent $\mathcal N$; it is unique and commutes with
base change. For a local base, the following are equivalent: $Q$ is locally
free; the displayed functor is right exact; projection formula is an
isomorphism for every $\mathcal N$; and base-change on $H^0$ is surjective.
Vanishing of $H^1(X_s,\mathcal F_s)$ implies these conditions.
\end{definition}

\begin{remark}[Universal functions]\label{kps-ex-universal-functions}
\source{kleiman-picard:page-0025}
If $f$ is proper and flat with reduced connected geometric fibers, then
$\mathcal O_S\risom f_*\mathcal O_X$ universally.
\end{remark}

\begin{definition}[Linear-system functor]\label{kps-def-linear-system-functor}
\dcref{3.12}\source{kleiman-picard:page-0025}
For invertible $\mathcal L$ on $X$, let $\LinSys_{\mathcal L/X/S}(T)$ be the
relative effective divisors $D$ on $X_T/T$ for which
\[\mathcal O_{X_T}(D)\risom\mathcal L_T\ox f_T^*\mathcal N\]
for some invertible $\mathcal N$ on $T$.
\end{definition}

\begin{theorem}[Linear systems are projective bundles]\label{kps-thm-linear-systems}
\dcref{3.13}\source{kleiman-picard:page-0025}
If $X/S$ is proper and flat with integral geometric fibers, and $Q$ is attached
to $\mathcal F=\mathcal L$, then $\PP(Q)$ represents
$\LinSys_{\mathcal L/X/S}$.
\uses{kps-def-Q,kps-def-linear-system-functor,kps-lem-fully-faithful,kps-lem-relative-criterion,kps-ex-universal-functions}
\end{theorem}
\begin{proof}
A divisor gives a unique (up to isomorphism) invertible $\mathcal N$ and a
regular section of $\mathcal L_T\ox f_T^*\mathcal N$. By the defining property
of $Q$, this is a map $u:Q_T\to\mathcal N$. Fiberwise nonvanishing and
Nakayama make $u$ surjective, hence a point of $\PP(Q)$. Conversely, the
tautological quotient yields a section whose restriction to every integral
fiber is nonzero, hence regular; the local criterion gives a relative effective
divisor. These constructions are inverse and functorial.
\end{proof}

\begin{remark}[Universal divisor]\label{kps-ex-universal-divisor}
\source{kleiman-picard:page-0026}
There is a universal relative divisor $W$ on $X_{\PP(Q)}$ with
$\mathcal O(W)\risom\mathcal L_{\PP(Q)}\ox f^*\mathcal O(1)$, characterized by
pullback along maps to $\PP(Q)$.
\end{remark}
